\section{Signed comparison obstructions and complete-space selection (v1.52)}
\label{sec:v152-route-selection}
The new material has three classifications: exact identities and proved
implications with stated hypotheses; two certified computations at an
already studied window; and explicit open inputs. No diagnostic is used
as a lower bound. A blocked parent route is not declared impossible
merely because one of its sufficient comparisons fails.

The direct uniform ordinary lower floor remains the weakest sufficient
cofinal target established here. CAE and unrestricted exact adaptive
concentration express that target in other coordinates. The additional
CCM problem is selection of a specific complete ground profile, with
entire-transform control. The following results remove specified
shortcuts and identify the estimates still missing; they do not supply
a new cofinal signed floor, G2 or RH.



\subsection*{Fresh certificates for two existing comparison failures}

The following computations use the existing $\lambda=5$ window and
the already stipulated metrics and partition. They introduce neither a
new window nor a replacement tail metric. The frozen vectors are new;
the original v1.27--v1.28 evidence is still not archived. In particular,
the historical displayed witness ratios and four-block tables are not
claimed to have been recovered or replayed.

\begin{proposition}[Fresh failure of the stipulated $10^{-8}D$ tail comparison]
\label{prop:v152-tail-comparison}
At $\lambda=5$, in each parity there is an exact dyadic vector supported
on $17\le n\le128$ with strictly positive Weil energy, but with
\[
 0<\frac{v^*\mathsf W_5v}{v^*Dv}
 <\begin{cases}7\cdot10^{-18},&\text{even},\\
               12\cdot10^{-18},&\text{odd}.
   \end{cases}
\]
Here $D=\operatorname{diag}(a_n)$ is the stipulated positive
archimedean metric, not the complete diagonal of $\mathsf W_5$.
Consequently $Q_{16}\mathsf W_5Q_{16}\succeq10^{-8}Q_{16}DQ_{16}$
is false in both parities. No negative Weil direction is asserted.
\end{proposition}
\begin{proof}[Certified verification]
The integer coordinates and common denominator $2^{60}$ are frozen in
\path{evidence/ns44_metric_replay/witnesses.json}. With those identical
vectors, fresh Arb coefficient enclosures at $320$ and $448$ bits,
including the analytic $96$-term series remainder, prove positive norm,
positive $D$ energy, positive Weil energy, and
$v^*(\mathsf W_5-10^{-8}D)v<0$ in each parity.
The respective ratio balls are contained in the displayed rational
intervals; rounded display values are
$6.4235244156530\cdot10^{-18}$ and
$1.1127922753304\cdot10^{-17}$.
The shifted energies are respectively approximately
$-2.4267798261306\cdot10^{-8}$ and
$-2.4549167289288\cdot10^{-8}$; the signs are decided by the
archived balls, not these rounded values.

A separate expansion into the normalized $+n,-n$ Fourier coordinates
verifies the same signs and overlaps the parity assembly. This checks
signs and indexing independently while sharing the analytic coefficient
enclosures; it is not a second derivation of the kernel. A further
reviewer replay and direct four-sign double sum reproduce the gates.
All coefficients are freshly assembled with an empty temporary cache.
The float proposal is used only to choose frozen vectors; no midpoint
solve enters a proof gate. The complete ball strings, source hashes,
verifier and independent review are archived and indexed in
\path{evidence/v152/}.
\end{proof}

\begin{proposition}[Fresh obstruction for the existing dyadic block comparison]
\label{prop:v152-block-comparison}
For the $\lambda=5$, $N=25$ dyadic partition, the signed block
norm-comparison criterion fails in both parities for every admissible
choice $0<M_j\preceq T_{jj}$ and every positive scalar row weighting.
The first three blocks
\[
 I_0=\{26,\ldots,50\},\quad I_1=\{51,\ldots,100\},\quad
 I_2=\{101,\ldots,200\}
\]
already suffice. This closes this partition's sufficient comparison,
not all signed block methods or the sign of the complete Weil form.
\end{proposition}
\begin{proof}[Certified pairings and exact comparison argument]
For each of the three pairs in each parity, separate frozen dyadic
vectors $x,y$ satisfy, at both $320$ and $448$ bits,
\[
 X=x^*T_{jj}x>0,\qquad Y=y^*T_{kk}y>0,\qquad
 |x^*T_{jk}y|^2-H_{jk}^2XY>0.
\]
The zero-diagonal symmetric rational matrices $H$ have entries
\[
 (H_{01},H_{02},H_{12})=
 \begin{cases}(997,568,481)/1000,&\text{even},\\
              (997,484,612)/1000,&\text{odd}.
 \end{cases}
\]
The same fresh coefficient and independent signed-index checks are
used for both energies and the cross pairing. The exact Rayleigh
quotients on $(1,1,1)$ are
\[
 \frac{341}{250}=1.364>1\quad\text{and}\quad
 \frac{2093}{1500}>1,
\]
respectively. The frozen witnesses and strict interval gates are in
\path{evidence/ns45_block_replay/}; an independent rerun gives a
byte-identical certificate. Pair witnesses need not be one common
physical vector because they lower-bound separate operator norms.

If admissible positive metrics exist, then each $T_{jj}$ is positive
definite. The verified quotients lower-bound the normalized coupling
norms for $T$, and Lemma~\ref{lem:v128-comparison-obstruction} gives
$\beta_{jk}(M)\ge\beta_{jk}(T)>H_{jk}$. If the metrics do not exist,
the criterion already fails its prerequisite. Thus no unconditional
full-block inertia assertion is needed. A positive row weighting making
every comparison row sum below one would, by diagonal similarity and
the row-sum norm, force the spectral radius of $H$ below one, contrary
to its exact Rayleigh quotient. An infinite continuation adds only
nonnegative norm terms, so it cannot repair these first three rows.
This proves precisely the stated comparison failure.
\end{proof}


\subsection*{Prime-channel retention for the explicit Gaussian weight}

% NS-43 working proof fragment. Exact identities and proved implications only.
% Not a manuscript/version change. No numerical certificate is asserted.

\paragraph{Notation and the comparison being ruled out.}
Write $C_0=(2\pi-3)/(4\pi)$ and retain the unnormalized Gaussian weight
$h$ of \eqref{eq:v150-weight}. For a prime power $n=p^k\ge2$, put
$\ell_n=\log n$, $w_n=\Lambda(n)/\sqrt n$, and
\[
 J^{(n)}_{a,h}[g]
 =w_n\int_{-a}^{a-\ell_n}h(x)h(x+\ell_n)
                 |g(x+\ell_n)-g(x)|^2\,dx
\]
when $\ell_n<2a$, and set it to zero otherwise. The global channel
$J^{(n)}_h$ has the same integral over $\mathbb R$. With $g=f/h$ on
$I_a$ and zero outside, $J^{(n)}_h\ge J^{(n)}_{a,h}$.

For $0\le\varepsilon_{n,a}\le1$, a coefficientwise loss is
$L_a[g]=\sum_n\varepsilon_{n,a}J^{(n)}_h[g]$ in the global comparison,
or the corresponding sum of interior channels in the local comparison.
Any additional discarded nonnegative energy can be included in $L_a$.
The two comparisons at issue are, respectively,
\begin{align*}
 J_h[g]-L_a[g]&\ge\tfrac23\operatorname{Var}_\nu(g)
       +2|\langle f,s\rangle|^2-\eta_a\|f\|^2,\\
 \int(r_a)_-|f|^2&\le J_{a,h}[g]-L_a[g]+\Pi[f]
                                  +\eta_a\|f\|^2.
\end{align*}
They require, respectively, $q_a[f]-L_a[g]\ge-\eta_a\|f\|^2$
and $q_a[f]-\int(r_a)_+|f|^2-L_a[g]\ge-\eta_a\|f\|^2$.
Thus both imply $\eta_a\ge L_a[g]-q_a[f]$ on every admissible unit test.
No compensating change in the potential, variance, pole or retained
coefficients is part of this statement.

\begin{proposition}[Every fixed nonzero prime-power channel has a critical coefficient]
\label{prop:ns43-fixed-channel}
Fix a prime power $n\ge2$, let $b=(\log2)/16$, and put
\[
 \kappa_n=C_0\Lambda(n)n^{-5},\qquad
 \gamma_n=\pi(1-n^{-2})e^{-6b}.
\]
Choose two fixed real even orthonormal functions
$\phi_1,\phi_2\in C_c^\infty(-b,b)$, and define
\[
 A_b=2\sum_{j=1}^2\int_{\mathbb R}
    |\Re\psi(1/4+i\xi/2)-\log\pi|\,|\widehat\phi_j(\xi)|^2\,d\xi,
 \qquad B_{a,b}=A_b+(8a+1)e^a+2a.
\]
For every sufficiently large $a$, each of the even actual-source
complement and the odd sector contains a unit smooth test satisfying
\[
 J^{(n)}_{a,h}[f/h]\ge\kappa_n\exp(\gamma_n e^{2a}),
 \qquad q_a[f]\le B_{a,b}.
\]
Consequently either comparison above, if valid on the indicated sector
with $\eta_a\ge0$, must obey
\[
 \eta_a\ge
 [\varepsilon_{n,a}\kappa_n\exp(\gamma_n e^{2a})-B_{a,b}]_+.
\]
In particular, deleting this channel, or decreasing its coefficient by
any fixed positive amount along a cofinal subfamily, is incompatible
with a uniform finite ordinary-norm error. This concerns the comparison,
not the sign of $q_a$ or the unmodified coefficient-one comparison.
\end{proposition}
\begin{proof}
The existing two-sided bounds for $h$ give, for $x\ge\ell_n$,
\[
 \frac{h(x-\ell_n)}{h(x)}
 \ge C_0 n^{-9/2}
             \exp\!\bigl(\pi(1-n^{-2})e^{2x}\bigr).
\]
The restriction $x\ge\ell_n$ is required here: both weight bounds
are being applied at nonnegative arguments.

Put $t=a-2b$ and
$F_{j,t}^{\pm}=2^{-1/2}(\phi_j(\cdot-t)\pm\phi_j(\cdot+t))$.
For $a>3b+\ell_n$, their support bands
$B_+\subset(a-3b,a-b)$ and $B_-=-B_+$ are disjoint and inside $I_a$.
The two $F_{j,t}^+$ are orthonormal; a unit vector in their complex
span can be chosen exactly orthogonal to the actual $u_a$, by the
kernel of the single linear functional $f\mapsto\langle f,u_a\rangle$.
This argument does not identify $u_a$ with $h$, approximate its
coefficients or require a source error estimate. The two odd functions
are orthonormal and every vector in their span is source-orthogonal.

For every unit vector in either span, an endpoint $x\in B_+$ has
$x-\ell_n\ge0$, remains inside the window, and is outside both support
bands, since $\ell_n\ge\log2>2b$. Select the edge from $x-\ell_n$ to
$x$, and use its reflected inward edge for $x\in B_-$. These selected
edges are distinct; no selected edge has two nonzero endpoints. Hence
\[
 J^{(n)}_{a,h}[f/h]
 \ge w_n\inf_{x\in(a-3b,a-b)}\frac{h(x-\ell_n)}{h(x)}
             \int_{B_+\cup B_-}|f(x)|^2dx
 \ge\kappa_n\exp\!\bigl(\pi(1-n^{-2})e^{2(a-3b)}\bigr).
\]
All interference between $\phi_1$ and $\phi_2$ is retained inside
$|f|^2$; its integral is the ordinary norm, exactly one.

The archimedean contribution is at most $A_b$: translation/reflection
multipliers have modulus at most $\sqrt2$, and Cauchy--Schwarz in the
two-dimensional coefficient vector gives the displayed bound uniformly
for every unit mixture. Prime correlations have total absolute value
at most $2\sum_{m<e^{2a}}\Lambda(m)/\sqrt m\le8ae^a$.
The even pole is at most $2a+2\sinh a\le2a+e^a$; the odd pole is
nonpositive. These estimates give $q_a[f]\le B_{a,b}$ including all
cross terms. The preceding comparison identities finish the proof.
Exterior edges were retained in those identities; the lower bound uses
only interior edges, so adding the exterior part cannot weaken it.
\end{proof}

\begin{proposition}[Uniform obstruction for moving omitted channels away from the endpoint]
\label{prop:ns43-moving-channel}
Fix $0<\theta<1$. Let
\[
 b_\theta=\frac{\min(\log2,-\log\theta)}{16},\qquad
 d_\theta=\min(\log2,-\log\theta-6b_\theta),
\]
\[
 \gamma_\theta=\pi(1-e^{-2d_\theta})e^{-6b_\theta}>0,
 \qquad c_\theta=C_0(\log2)\theta^{-1/2}e^{9b_\theta/2}>0.
\]
For all sufficiently large $a$, the same two parity spaces built with
$b=b_\theta$, $t=a-2b_\theta$, satisfy, for every unit vector and
\emph{simultaneously for every prime power} $2\le n\le\theta e^{2a}$,
\[
 J^{(n)}_{a,h}[f/h]\ge
 c_\theta e^{-11a/2}\exp(\gamma_\theta e^{2a}).
\]
Thus either comparison, with $\eta_a\ge0$, and an omitted channel
$n=n(a)$ in this range, even when it varies with $a$, forces
\[
 \eta_a\ge
 [\varepsilon_{n(a),a}c_\theta e^{-11a/2}
             \exp(\gamma_\theta e^{2a})-B_{a,b_\theta}]_+.
\]
The same statement holds on any cofinal subfamily satisfying these
hypotheses. No conclusion is asserted merely because an omitted
channel exists at each window if its ratio $n(a)/e^{2a}$ tends to one.
\end{proposition}
\begin{proof}
Here $d_\theta\ge10b_\theta>2b_\theta$. For $x\in(a-3b,a-b)$ and
$\log2\le\ell_n\le2a+\log\theta$, one has
\[
 \ell_n\ge d_\theta,\qquad 2x-\ell_n\ge d_\theta,
 \quad\text{hence}\quad |x-\ell_n|\le x-d_\theta.
\]
For large $a$ these inward endpoints lie inside the window and outside
both support bands, since
$|x-\ell_n|\le a-b-d_\theta<a-3b$. Unlike the fixed-channel proof,
$x-\ell_n$ need not be nonnegative. Use evenness of $h$ and apply its
bounds to $|x-\ell_n|$ and $x$, rather than incorrectly applying the
previous ratio formula past its range. They give
\begin{align*}
 \frac{h(x-\ell_n)}{h(x)}
 &\ge C_0\exp\!\left(\tfrac92(|x-\ell_n|-x)
                   +\pi(e^{2x}-e^{2|x-\ell_n|})\right)\\
 &\ge C_0e^{-9(a-b)/2}
                \exp\!\left(\pi(1-e^{-2d_\theta})e^{2(a-3b)}\right).
\end{align*}
Also $w_n\ge(\log2)\theta^{-1/2}e^{-a}$. The same distinct selected
inward edges and norm-one integration prove the claimed bound. The
source kernel and the original-form upper budget are exactly as above,
with fixed $b_\theta$ independent of $a$ and $n$.
\end{proof}

\paragraph{Truncations actually covered.}
With every retained channel coefficient at most one and all other
signed terms unchanged, the following shortcuts require unbounded error:
(i) any fixed finite set of retained prime-power channels; (ii) any
fixed finite set of retained primes, even retaining all of their powers;
(iii) a uniformly bounded number of retained channels (or primes),
although their labels may vary with the window; and (iv) a cutoff
retaining only $n\le M(a)$ with $M(a)=o(e^{2a})$.
For (iii), take more fixed prime powers (respectively primes) than the
retained count: one is omitted, and all lie in the uniform range for
large $a$. For (iv), if $M(a)\ge1$, the first power of $2$ exceeding
$M(a)$ is at most $2M(a)$ and is eventually at most $\theta e^{2a}$;
if $M(a)<1$, take $n=2$ instead. The same proof covers
$\limsup M(a)e^{-2a}<1/2$ by choosing $\theta$ between twice that
limsup and one. No prime-gap theorem is being used.

A fixed coefficient loss in any of those omitted channels is covered.
A loss tending sufficiently rapidly to zero is only subject to the
necessary displayed rate; it is not excluded categorically. Merely
saying ``finitely many channels at each window'' is insufficient:
the full local sum itself is finite. A cutoff approaching the strict
endpoint, compensating changes in the signed diagonal/pole, coefficients
larger than one in other channels, and a different positive weight are
not ruled out by these propositions.

\paragraph{One specific surviving signed, cofinal next input.}
This is a proposed \emph{open input on a restricted test family}, not
a theorem and not a restatement of CAE on all tests. For fixed real even
$\phi_1,\phi_2\in C_c^\infty(-b,b)$ as above, set
\[
 R_{jk}(u)=\int\phi_j(z+u)\phi_k(z)\,dz,\qquad
 C_j=\int\cosh(x/2)\phi_j(x)\,dx,
\]
and consider the explicit real symmetric two-by-two matrix
\[
 E_{jk}(t)=e^t C_jC_k-
       \sum_{n\ge2}\frac{\Lambda(n)}{\sqrt n}
                          R_{jk}(\log n-2t),\qquad t>2b.
\]
Only $e^{2t-2b}<n<e^{2t+2b}$ contribute. The candidate next lemma is
\[
 \sup_{t\ge t_0}\|E(t)\|_{\mathrm{op}}<\infty
 \quad\text{for this specified pair of bumps.}
\]
This asks for a signed, smoothly weighted prime--continuum cancellation
at cofinal multiplicative locations, not separate absolute estimates
of Gaussian transformed channels.

Its limited consequence is checkable. Let
$D_{jk}=q(\phi_j,\phi_k)$ and
$C_{jk}(t)=q(\phi_j(\cdot-t),\phi_k(\cdot+t))$ for the polarized
physical form. For $a=t+2b$ the complete prime and pole formula gives
\[
 C_{jk}(t)=E_{jk}(t)+e^{-t}C_jC_k+A_{jk}(t),
 \qquad A_{jk}(t)=O_b(e^{-t}),
\]
where the last term is the off-diagonal archimedean form. Indeed its
kernel on disjoint supports is $-\rho(|x-y|)$, so the decay follows
from $\rho(s)=O(e^{-s/2})$. The pole is
$2C_jC_k\cosh t=(e^t+e^{-t})C_jC_k$; precisely one orientation of
the prime shift connects the separated supports, giving the sum above
without an extra factor two. The parity matrices are exactly
$D+C(t)$ and $D-C(t)$, with all cross terms present. The candidate
bound therefore supplies a uniform floor only on these two-dimensional
parity spaces, including their actual-source-orthogonal even vectors.
It does not cover arbitrary physical tests, does not establish CAE,
and is not asserted to be weaker than RH or presently provable.


\subsection*{A restricted prime-discrepancy input can already carry the full zero obstruction}
\label{sec:ns43-bump-strength}
The two-dimensional test family in the channel analysis is concrete,
but its uniform discrepancy estimate must not be advertised as a
weaker arithmetic foothold without checking its strength.

\begin{proposition}[What bounded translated-bump discrepancy excludes]
\label{prop:ns43-bump-strength}
Let $\phi\in C_c^\infty(-b,b)$ be real and even, put
\[
 P(s)=\int e^{sx}\phi(x)\,dx,\qquad
 R(v)=\int\phi(x+v)\phi(x)\,dx,\qquad C=P(1/2),
\]
and define
\[
 F(u)=e^{u/2}C^2-
        \sum_{n\ge2}\frac{\Lambda(n)}{\sqrt n}R(\log n-u).
\]
If $F$ is bounded for $u\ge u_0$, every nontrivial zero $\rho$ of
$\zeta$ with $\Re\rho>1/2$ must obey $P(\rho-1/2)=0$.
There are explicit nonnegative even smooth compactly supported probes
$\phi$ for which all zeros of $P$ lie on the imaginary axis. For
such a probe, the asserted boundedness would exclude every zero off
the critical line. This is a conditional implication, not a proved
bound for $F$ and not a new RH claim. For an arbitrary preassigned
pair of bumps, the matrix input forces common transform zeros at
any off-line zero; no conclusion without that qualification is asserted.
\end{proposition}
\begin{proof}
The correlation is supported in $[-2b,2b]$ and satisfies
$\int e^{sv}R(v)\,dv=P(s)^2$ by evenness. For $\Re s>1/2$,
absolute convergence of the von Mangoldt Dirichlet series gives
\begin{equation}\label{eq:ns43-bump-laplace}
 \int_{u_0}^{\infty}e^{-su}F(u)\,du
 =\frac{C^2}{s-1/2}+P(s)^2\frac{\zeta'}{\zeta}(s+1/2)+H(s),
\end{equation}
where
\begin{align*}
 H(s)={}&C^2\frac{e^{-(s-1/2)u_0}-1}{s-1/2}\\
 &+\sum_{\log n\le u_0+2b}\frac{\Lambda(n)}{\sqrt n}
       \int_{-\infty}^{u_0}e^{-su}R(\log n-u)\,du.
\end{align*}
This correction is entire: the quotient is removable and only
finitely many compactly supported integrals occur. The usual
identity $-\zeta'/\zeta(s)=\sum\Lambda(n)n^{-s}$ is used only in
its half-plane of absolute convergence (\href{https://dlmf.nist.gov/27.4.E12}{DLMF 27.4.12}). The pole at $s=1/2$ in
\eqref{eq:ns43-bump-laplace} cancels because $C=P(1/2)$.
Boundedness of $F$ makes the left side holomorphic on $\Re s>0$.
Uniqueness of meromorphic continuation therefore forces the residue
at $s=\rho-1/2$ to vanish. That residue is
$\operatorname{mult}(\rho)P(\rho-1/2)^2$, proving the first assertion.
For a matrix of two bumps apply this argument to each diagonal;
all off-line zeros must then be common zeros of their transforms.

Here is a zero-independent choice of probe. Fix $0<B<b$ and set
$b_k=B/(\zeta(2)k^2)$. Let $\mu$ be the distribution of
$\sum_{k\ge1}U_k$, where the $U_k$ are independent uniform variables
on $[-b_k,b_k]$. The sum converges absolutely since $\sum b_k=B$.
Its measure is nonnegative, even, supported in $[-B,B]$, and has
entire bilateral transform
\[
 P_0(s)=\prod_{k\ge1}\frac{\sinh(b_ks)}{b_ks}.
\]
The product converges uniformly on compact sets, since each factor
is $1+O(b_k^2)$ there and $\sum b_k^2<\infty$. It is nonzero away
from the zeros of its factors. Those zeros are exactly imaginary
points $i\pi m/b_k$, $m\ne0$.
For every integer $M>0$,
\[
 |P_0(i\xi)|\le
 \left(\prod_{k=1}^M b_k^{-1}\right)|\xi|^{-M},
\]
since all remaining sinc factors have modulus at most one. Thus
Fourier inversion gives a $C^\infty$ density supported in $[-B,B]$.
It is nonnegative and nonzero because it is the density of $\mu$;
normalizing it in $L^2$ does not change its transform zeros. This is
the required $\phi\in C_c^\infty(-b,b)$. If an orthonormal pair is
desired, take the second even smooth bump supported in
$\{B<|x|<b\}$ and normalize it; disjoint supports give orthogonality.
For the first probe no zero with $\Re(\rho-1/2)>0$ can be masked.
The functional-equation symmetry supplies the corresponding exclusion
on the other side of the critical line, conditional on the unproved
boundedness of $F$.
\end{proof}

For the proposed matrix in the channel analysis, $F(u)=E_{11}(u/2)$
for this first probe. Hence even this restricted uniform estimate can
already contain the full off-line-zero obstruction. Nothing above
proves that estimate or asserts the converse. The useful outcome of
this check is a limit on the advertised strength of a proposed next
lemma, not a new positivity argument.


% NS-43 concentration subtask. Analytic working fragment, not a version claim.
% All conclusions are explicitly scoped. No numerical or RH claim.
\subsection*{Adaptive concentration: completeness is not an estimate (NS-43)}
\label{ns43-conc-sec}

\paragraph{Objects and source constraint.}
Fix $a>0$. Let $\mathcal H_a=L^2(-a,a)$, let $\mathcal F_a$ mean
zero extension followed by the unitary Fourier transform, and let
$\beta_a$ be the exact complete symbol, including its endpoint and pole
terms. Put $D_a=\| (\beta_a)_-\|_\infty$.
Use either of the closed spaces
\[
 S_a^{\rm ev}=\{f\in L^2_{\rm ev}(-a,a):\langle f,u_a\rangle=0\},
 \qquad S_a^{\rm odd}=L^2_{\rm odd}(-a,a),
\]
where $u_a$ is the actual normalized repaired source.
In this subsection $S_a$ denotes one of these spaces and $P_{S_a}$
its orthogonal projection. All operator lower bounds act on the
\emph{whole} $S_a$, not on a head or on individual pencil directions.
For an even Borel set $G$ put
\[
 C_{a,S}(G)=P_{S_a}\mathcal F_a^*1_G\mathcal F_aP_{S_a}|_{S_a}.
\]
The complete form is
$q_a[f]=\int\beta_a|\mathcal F_af|^2$, with form domain given by
$\int(\beta_a+D_a)|\mathcal F_af|^2<\infty$.

\begin{proposition}[Exact adaptive hierarchy and the original lower edge]
\label{ns43-conc-prop-adaptive}
Suppose $\beta_a$ is continuous, even, bounded below, and tends to
$+\infty$ at both frequency ends. Assume also that
$S_a\cap\mathcal D(q_a)$ contains a nonzero vector.
These hypotheses hold for the manuscript's fixed-window symbol:
$\beta_a(\xi)=O_a(\log(2+|\xi|))$ puts every compactly supported
smooth physical test in the form domain. In the even sector, a
nonzero linear combination of two independent even smooth tests
can satisfy the one source constraint; in the odd sector any
nonzero odd smooth test suffices. For integers $n\ge1$ define
\begin{align*}
 \delta_n&=2^{-n},\qquad J_n=n2^n,\\
 m_{a,n}(\xi)&=-D_a+
 \delta_n\sum_{j=1}^{J_n}
       1_{\{\beta_a(\xi)+D_a\ge j\delta_n\}}.
\end{align*}
These are finite step minorants with positive weights, and
$m_{a,n}\uparrow\beta_a$ pointwise. Their exact compressed lower
edges obey
\begin{equation}\label{ns43-conc-eq-edge-limit}
 \lambda_{a,n}:=\inf_{\substack{f\in S_a\\\|f\|=1}}
   \int m_{a,n}|\mathcal F_af|^2
 \ \uparrow\
 \lambda_a:=\inf_{\substack{f\in S_a\cap\mathcal D(q_a)\\\|f\|=1}}
   q_a[f].
\end{equation}
Consequently the following two existence statements on any prescribed
cofinal family are equivalent:
\begin{enumerate}[label=(\alph*)]
\item the complete even source complement and complete odd sector
have one common finite lower floor $-C$, independent of $a$;
\item one may choose finite indices $n_{\rm ev}(a),n_{\rm odd}(a)$
so that the corresponding exact signed concentration sums
\[
 -D_a I_{S_a}+\delta_n\sum_{j=1}^{J_n}
 C_{a,S}\bigl(\{\beta_a+D_a\ge j\delta_n\}\bigr)
\]
have a common finite lower floor independent of $a$.
\end{enumerate}
More precisely, (a) yields a floor $-C-\varepsilon$ in (b) for any
fixed $\varepsilon>0$; (b) yields its stated floor in (a).
This equivalence supplies no effective choice of $n(a)$ and no
arithmetic lower estimate for any of the sums.
\end{proposition}
\begin{proof}
The formula is lower dyadic rounding of
$\min(\beta_a+D_a,n)$, minus $D_a$. Refinement of the grid and
increase of the cap prove monotonicity and pointwise convergence.
Each $m_{a,n}$ is bounded. Thus its compressed form is a bounded
self-adjoint operator, and $\lambda_{a,n}\le\lambda_a$.
The numbers increase to some finite $\ell\le\lambda_a$; finiteness
above follows by testing one fixed unit vector in
$S_a\cap\mathcal D(q_a)$, and below by $-D_a$.

Choose unit $f_n\in S_a$ with
$\int m_{a,n}|\mathcal F_af_n|^2\le\lambda_{a,n}+1/n$.
The integrals against $m_{a,n}+D_a$ are uniformly bounded by some
$B<\infty$. Given $R\ge1$, choose $M_R$ with
$\beta_a(\xi)+D_a\ge2R$ for $|\xi|>M_R$.
For sufficiently large $n$, lower rounding gives
$m_{a,n}+D_a\ge R$ there. Hence
\[
 \int_{|\xi|>M_R}|\mathcal F_af_n|^2\le B/R.
\]
Restriction of $\mathcal F_a$ to $[-M_R,M_R]$ is a compact map:
its integral kernel has finite square integral on the product of
the two bounded intervals. This fact and the uniform frequency-tail
bound make $(\mathcal F_af_n)$, after discarding finitely many
terms for each $R$, relatively compact in $L^2(\mathbb R)$.
Take a strongly convergent subsequence. Its physical limit $f$ is
in the closed space $S_a$ and has norm one.

For each fixed $k$, boundedness of $m_{a,k}$ and monotonicity give
\[
 \int m_{a,k}|\mathcal F_af|^2
 =\lim_n\int m_{a,k}|\mathcal F_af_n|^2\le\ell.
\]
Monotone convergence applied after adding $D_a$ gives
$f\in\mathcal D(q_a)$ and $q_a[f]\le\ell$.
Thus $\lambda_a\le\ell$, proving \eqref{ns43-conc-eq-edge-limit}.
Choose the two finite indices separately at each window to within
$\varepsilon$ of their limits. Conversely a pointwise minorant's
operator lower bound transfers immediately to $q_a$.
\end{proof}

\paragraph{What this closes.}
The proposal ``allow arbitrary adaptive levels and assert that their
exact signed operator sums have a uniform lower floor'' is an
equivalent restatement of the unknown complement-floor target.
It is not a weaker arithmetic input. This closes only that claimed
logical shortcut. It does not close a method that proves a specific
operator comparison from additional arithmetic structure.
The finite level count is allowed to grow without restriction here.
Therefore the statement does not evade a proved obstruction by
silently retaining bounded complexity.

\begin{proposition}[All scalar set bounds can still lose the signed floor]
\label{ns43-conc-prop-scalar}
Even exact upper and lower concentration bounds for every measurable
set do not in general determine the signed operator lower edge.
This failure occurs already for commuting positive operator-valued
measures with three outcomes, after an exact rank-one source has
been removed.
\end{proposition}
\begin{proof}
On $\mathbb C^3$ define two three-outcome positive operator-valued
measures by making outcome operator $C_i$ diagonal, with its three
diagonal entries the $i$th column of the following arrays:
\[
 A=\frac1{10}
 \begin{pmatrix}0&7&3\\3&0&7\\7&3&0\end{pmatrix},
 \qquad
 B=\frac1{10}
 \begin{pmatrix}7&0&3\\3&7&0\\0&3&7\end{pmatrix}.
\]
Each row sums to one, so in both models $\sum_iC_i=I$.
For every singleton $E$, the exact lower and upper scalar bounds are
$p(E)=0$, $c(E)=7/10$. For every two-element set they are
$p(E)=3/10$, $c(E)=1$. Empty and full sets have their usual values.
These list \emph{all} subsets; hence all scalar data coincide.

For symbol values $b=(-1/2,1/2,3/2)$, the exact signed operators have
diagonal entries
\[
 \sum_i b_i C_i^A=\operatorname{diag}(4/5,9/10,-1/5),
 \qquad
 \sum_i b_i C_i^B=\operatorname{diag}(1/10,1/5,6/5).
\]
Thus their bottoms are $-1/5$ and $1/10$ despite identical scalar
concentration data for every set.

The usual scalar certification rule cannot recover the positive
bottom in model $B$, even after optimizing over all sets and weights.
Indeed the probability vector $z=(7/10,3/10,0)$ satisfies every
scalar constraint $p(E)\le z(E)\le c(E)$ in model $B$.
For any pointwise minorant
$\alpha+\sum_jw_j1_{G_j}\le b$ with $w_j\ge0$, one therefore has
\[
 \alpha+\sum_jw_jp(G_j)
 \le\alpha+\sum_jw_jz(G_j)\le\sum_i b_i z_i=-1/5.
\]
This upper bound on the certifiable lower bound is attained by
$b=-1/2+1_{\{2,3\}}+1_{\{3\}}$.
Multiplying $b$ by any $t>0$ gives an exact best scalar lower bound
$-t/5$, while the actual model-$B$ operator is at least $(t/10)I$.
The scalar loss therefore diverges along $t\to\infty$ even though
the true signed floor is nonnegative.

To add an exact source, take an orthogonal direct sum with one unit
vector $u$ and give the three outcome operators on that vector the
values $(1/2,1/2,0)$. Their signed sum annihilates $u$ and has no
cross terms with $u^\perp$. Compression off this same exact source
returns the two arrays above. The countermodel is finite and
nonarithmetic: it does not assert that either array arises from
physical Paley--Wiener concentration or from the corrected zero field.
\end{proof}

\paragraph{Relation to the existing obstructions.}
The negative-only and primitive scalar estimates discard signed
physical information; v1.37 closes its stated pointwise and primitive
uniform-budget proposals. The signed operators in
Proposition~\ref{ns43-conc-prop-adaptive} retain the positive regions
and all cross terms, so those scalar no-go proofs do not apply to them.
This is a distinction in hypotheses, not a proof that their lower
edges meet the target.

The QG growth clause in Assumption~\ref{ass:v149-qg}, together with
its \emph{separate} source-admissible unit packet satisfying the
all-Borel domination and uniform upper energy assumptions, gives
\[
 \eta(m_a)\ge
 [c_0\,\mathrm{QC}_{K(a)}(\widetilde m_a)-Q_*]_+.
\]
It obstructs a bounded level count under those two open inputs.
For every fixed bound $K_0$, the cost diverges, so a successful
bounded-loss family would have $K(a)\to\infty$.
Neither QG alone, a finite-bin packet test, nor the conditional
obstruction proves failure of unrestricted adaptive comparison.
In particular there is no contradiction with
Proposition~\ref{ns43-conc-prop-adaptive}: its indices are unrestricted
and its existence is conditional on the original floor.

The odd sector requires no source projection. In the even sector,
the actual source must be removed before every concentration or
packet comparison. For a raw unit $f$ with
$r=|\langle f,u_a\rangle|<1$, its normalized projection $g$ obeys
the already proved pointwise and energy estimates
\[
 |\mathcal F_ag(\xi)|\ge
 \frac{[|\mathcal F_af(\xi)|-r|\mathcal F_au_a(\xi)|]_+}
 {\sqrt{1-r^2}},\qquad
 |(1-r^2)q_a[g]-q_a[f]|
 \le(2r+r^2)\|\mathsf W_{e^a}u_a\|.
\]
Thus raw packet coverage or a raw pencil ratio cannot be substituted
for the source-compressed hypotheses. The repaired-source residual
is used only in the subsequent complement-to-full-floor transfer;
it does not supply the missing signed comparison or packet coverage.

\paragraph{Disposition and missing input.}
The full signed concentration route remains open. The proposed
unrestricted exact hierarchy is classified as an \emph{exact
reformulation}; completeness of optimization from individual scalar
set bounds is a \emph{proved-false general simplification}, with no
arithmetic impossibility claim. A productive next input must estimate
joint signed operators using actual arithmetic structure, with a
cofinal uniform error on both complete parity obligations. No such
estimate is proved here. QG is an obstruction-investigation input,
not a positivity mechanism; CAE is already the same floor in another
coordinate system. These statements add no numerical window, tail
metric, certificate, G2 conclusion, or RH conclusion.


% NS-43 analytic working note. No numerical computation or certificate.
% Scope: additional complete-space inputs for CCM step (b).

\subsection*{CCM selection requires a relative estimate, not another absolute residual}

\paragraph{Status.}
The fixed-window simple ground and its real-zero transform do not identify
Riemann's $\Xi$. The known source limit and complete residual are useful
prerequisites. The new input would be a cofinal \emph{complete spectral
selection estimate}, with a rate compatible with the transform topology.
The propositions below are conditional elementary functional analysis;
none establishes the hypotheses for the arithmetic Weil family.

Put $H=L^2(\mathbb R)$, $I_a=(-a,a)$ and let $J_a$ denote zero extension.
Let $A_a$ be the complete lower-bounded self-adjoint even-sector operator
on $L^2(I_a)$, with compact resolvent. If the global ground is used,
its evenness and ordering below the odd sector must also be supplied.
In the following selection lemmas assume the selected ground is simple,
with unit vector $v_a$, eigenvalue $\mu_a$, and complete next-level
separation $g_a>0$:
\[
 q_a[f]-\mu_a\|f\|^2\ge g_a\|f\|^2
 \quad(f\in\mathcal D(q_a),\ f\perp v_a).
\]
This is an ordinary complete spectral separation, not an LDL pivot,
finite-compression gap, or a source-complement estimate inferred from
separate diagonal blocks.

Let $h$ be the explicit positive Gaussian arithmetic radical of
\eqref{eq:v150-weight}, and $\phi=h/\|h\|_2$. Use exactly the fixed
boundary profile of \ref{prop:v135-growing-radical}: fix one smooth
$\vartheta$ with $0\le\vartheta\le1$, $\vartheta(t)=1$ for $t\le-1$,
and $\vartheta(t)=0$ for $t\ge-1/2$. For $a\ge2$ set
$\chi_a(x)=\vartheta(|x|-a)$. This is smooth (it is constant near zero),
even, supported in $I_a$, equal to one on $[-a+1,a-1]$, and has
uniformly bounded derivatives of every fixed order. Define
\[
 p_a=\chi_a\phi/\|\chi_a\phi\|_2,\qquad
 \alpha_a=q_a[p_a],\qquad E_a=\alpha_a-\mu_a\ge0.
\]
For large $a$, these vectors lie in the complete operator domain.
The double-exponential tail from v1.35 implies, for every fixed $R\ge0$,
\[
 \int_{\mathbb R}e^{R|x|}|J_ap_a(x)-\phi(x)|\,dx\longrightarrow0.
\]
The vector $p_a$ is the normalized $j=0$ column of the growing block.
Thus \ref{prop:v135-growing-radical}, specifically
\eqref{eq:v135-block-residual}, gives the complete estimate
\[
 \|A_ap_a\|\le\varepsilon_a:=Ca\exp(-e^a/100),\qquad
 |\alpha_a|\le\varepsilon_a,
\]
with constants for this fixed source and profile. No residual estimate
is asserted for arbitrary cutoffs whose derivatives may grow with $a$.
The cutoff $p_a$ here is explicit and is not silently identified with
the different repaired prolate source $u_a$.

\begin{proposition}[A complete relative selection estimate]
\label{prop:ns43-ccm-rayleigh-selection}
Choose the phase of $v_a$ so that $\langle p_a,v_a\rangle\ge0$.
If $E_a/g_a\to0$, then $\|J_av_a-\phi\|_2\to0$.
Define
\[
 C_R(a)=\begin{cases}
 \sqrt{(e^{2Ra}-1)/(2\pi R)},&R>0,\\
 \sqrt{a/\pi},&R=0.
 \end{cases}
\]
If, in addition, for every fixed $R>0$,
\[
 C_R(a)\sqrt{E_a/g_a}\longrightarrow0,
 \tag{S1}
\]
then the unitary Fourier transforms of $J_av_a$ converge locally
uniformly on $\mathbb C$ to $\widehat\phi$.
\end{proposition}
\begin{proof}
Write $p_a=\kappa_a v_a+e_a$, where $e_a\perp v_a$ and
$\kappa_a\ge0$. The spectral projection preserves the form domain,
and the cross term is exactly zero because $v_a$ is an eigenvector.
Consequently
\[
 E_a=q_a[e_a]-\mu_a\|e_a\|^2\ge g_a\|e_a\|^2,
 \quad \kappa_a^2=1-\|e_a\|^2,
 \quad \|p_a-v_a\|^2=2(1-\kappa_a)\le2E_a/g_a.
\]
This proves the ordinary-norm conclusion. For $|\Im z|\le R$,
Cauchy--Schwarz gives
\[
 |\widehat{J_a(v_a-p_a)}(z)|
 \le (2\pi)^{-1/2}\left(\int_{-a}^a e^{2R|x|}\,dx\right)^{1/2}
       \|v_a-p_a\|
 \le C_R(a)\sqrt{2E_a/g_a}.
\]
The weighted cutoff convergence above controls
$\widehat{J_ap_a}-\widehat\phi$ on the same strip. Combining the
bounds proves the assertion. No compactness hypothesis is hidden:
(S1) supplies a direct quantitative topology upgrade.
\end{proof}

\begin{proposition}[A residual-to-separator version]
\label{prop:ns43-ccm-residual-selection}
Let $p_a\in\mathcal D(A_a)$ be unit, put
$\alpha_a=\langle A_ap_a,p_a\rangle$ and
$\rho_a=\|(A_a-\alpha_a)p_a\|$. Suppose an independent complete
second-eigenvalue lower bound $\beta_a>\alpha_a$ is established,
with eigenvalues counted with multiplicity. Then the ground is simple and
\[
 \|v_a-p_a\|\le\frac{\sqrt2\,\rho_a}{\beta_a-\alpha_a}.
\]
Thus the cofinal condition
\[
 C_R(a)\frac{\rho_a}{\beta_a-\alpha_a}\longrightarrow0
 \quad\text{for every fixed }R>0 \tag{S1r}
\]
selects the entire transform of the explicit cutoff radical.
\end{proposition}
\begin{proof}
The variational upper bound is $\mu_a\le\alpha_a<\beta_a$, so
exactly one eigenvalue lies below $\beta_a$ and it is simple.
On its orthogonal complement the spectrum of $A_a-\alpha_a$ lies
above $\beta_a-\alpha_a$. Spectral expansion gives
$\rho_a^2\ge(\beta_a-\alpha_a)^2\|e_a\|^2$.
The preceding phase/norm identity and Fourier estimate complete the proof.
\end{proof}
This formulation does not require a lower ground-energy estimate matching
$\alpha_a$ to high relative precision. The missing arithmetic input is
the complete cofinal separator at the correct relative scale. A finite
matrix's second eigenvalue supplies no such lower bound without its full
tail and coupling comparison. An even-sector separator still needs the
separate odd ordering if the real-zero theorem requires the global ground.

\paragraph{A weaker alternative to the evaluator rate.}
Ordinary selection $E_a/g_a\to0$ also gives locally uniform entire
convergence if the selected grounds satisfy the independently proved
weighted tightness condition
\[
 \lim_{K\to\infty}\sup_a\int_{|x|>K}e^{R|x|}|J_av_a(x)|\,dx=0
 \quad\text{for every }R>0. \tag{S2}
\]
Indeed the integral on $[-K,K]$ converges by ordinary $L^2$ convergence
and Cauchy--Schwarz, while (S2) and the known tail of $\phi$ control
the exterior. A sufficient explicit condition for (S2) is
$\sup_a\int e^{2(R+1)|x|}|J_av_a|^2<\infty$ for each $R$.
No such ground-tail statement follows from the current residual bound.

\begin{proposition}[Complete evaluator-dual alternative]
\label{prop:ns43-ccm-evaluator-selection}
In the preceding setting suppose $E_a/g_a\to0$. Suppose further that,
for every fixed $R>0$, a complete complement estimate is established:
\[
 \sup_{|z|\le R}|\widehat{J_ae}(z)|^2
 \le D_{a,R}\bigl(q_a[e]-\mu_a\|e\|^2\bigr)
 \quad(e\in\mathcal D(q_a),\ e\perp v_a),
 \qquad D_{a,R}E_a\to0. \tag{S3}
\]
Then the same locally uniform entire convergence follows.
\end{proposition}
\begin{proof}
For the orthogonal error above,
$\widehat p_a=\kappa_a\widehat v_a+\widehat e_a$ and
$\kappa_a\to1$. Thus, on the disk,
\[
 |\widehat v_a-\widehat p_a|
 \le \kappa_a^{-1}\sqrt{D_{a,R}E_a}
       +(\kappa_a^{-1}-1)|\widehat p_a|\to0.
\]
The explicit cutoff transforms are uniformly bounded on each disk
and converge to $\widehat\phi$. This proves the assertion.
\end{proof}
The shift by $\mu_a$ and the complete complement in (S3) are part of
its hypotheses. The fixed-window v1.43 evaluator comparison for one
real zero is not a cofinal bound of this kind, nor a bound on every
complex compact set. Its scalar resolution floor is not a no-go
result for (S1)--(S3).

\begin{proposition}[A source-relative block condition sufficient for selection]
\label{prop:ns43-ccm-source-separation}
Let $p_a\in\mathcal D(A_a)$ be unit, $P_a=|p_a\rangle\langle p_a|$,
$Q_a=I-P_a$, $\alpha_a=\langle A_ap_a,p_a\rangle$, and
$r_a=Q_aA_ap_a$. Suppose the \emph{complete compressed form} satisfies
\[
 q_a[f]\ge(\alpha_a+\delta_a)\|f\|^2
 \quad(f\in\mathcal D(q_a)\cap p_a^\perp),\qquad \delta_a>0.
 \tag{S4}
\]
Then $A_a$ has exactly one eigenvalue below $\alpha_a+\delta_a$,
it is simple, and, for its phase-aligned ground $v_a$,
\[
 \alpha_a-\|r_a\|^2/\delta_a\le\mu_a\le\alpha_a,
 \qquad \|v_a-p_a\|\le\sqrt2\,\|r_a\|/\delta_a.
\]
For the explicit cutoff source, the cofinal rate
$C_R(a)\|r_a\|/\delta_a\to0$ for every $R>0$ therefore implies
entire-transform selection.
\end{proposition}
\begin{proof}
The restriction to $p_a^\perp$ is a closed lower-bounded form,
because subtracting its component along the fixed operator-domain
vector preserves the form domain continuously. The min--max principle
places the second eigenvalue at least $\alpha_a+\delta_a$; the trial
$p_a$ puts the lowest at most $\alpha_a$. For $f=tp_a+y$, $y\perp p_a$,
\[
 q_a[f]-\alpha_a\|f\|^2
 =q_a[y]-\alpha_a\|y\|^2+2\Re(t\langle A_ap_a,y\rangle)
 \ge\delta_a\|y\|^2-2|t|\|r_a\|\|y\|
 \ge-\|r_a\|^2|t|^2/\delta_a.
\]
This handles the full cross term and proves the lower edge bound.
Writing $v_a=tp_a+y$, the weak projected eigen-equation is
$(Q_aA_aQ_a-\mu_a)y=-tr_a$. Its inverse has norm at most
$1/(\alpha_a+\delta_a-\mu_a)\le1/\delta_a$, so
$\|y\|\le|t|\|r_a\|/\delta_a$. Phase alignment then gives
$\|v_a-p_a\|^2=2(1-|t|)\le2\|y\|^2$.
\end{proof}
Condition (S4) is a new, source-specific signed spectral estimate,
not an algebraic reformulation of the bounded-floor target. It asks for information beyond a floor: it separates one named source
from its whole complement and controls the relative error. No reverse
implication from RH or from a uniform floor to (S4) is claimed. In this
arithmetic family, with the known vanishing source residual, a successful
cofinal condition of this strength gives a uniform floor in the sector
where it is imposed. To use v1.36 one must impose it on the full space
or separately control the odd sector by a uniform floor; those additional
hypotheses would imply RH. Even-sector selection does not silently supply
the odd estimate. It is not an unconditional shortcut.

For clarity, the floor consequence itself needs no relative-rate
condition: in the displayed cross-term formula,
$2|t|\|y\|\le |t|^2+\|y\|^2$ and $\delta_a>0$ give
\[
 q_a[tp_a+y]\ge(\alpha_a-\|r_a\|)(|t|^2+\|y\|^2).
\]
Thus the known $\alpha_a\to0$ and $\|r_a\|\to0$ give the asserted
sector floor from (S4) alone. The stronger relative-rate condition
is needed for selection of the Gaussian profile and its transform,
not for this floor deduction.

The v1.35 even reflected near-radical block has dimension $J_a+1$,
so eventually it contains a unit vector orthogonal to $p_a$. Its complete
residual is at most $\varepsilon_a$, hence (S4) forces
$\delta_a\le\varepsilon_a-\alpha_a\le\varepsilon_a+|\alpha_a|$.
The explicit cutoff source also has $|\alpha_a|=O(\varepsilon_a)$.
This rules out a uniform positive $\delta_a$ or a polynomial lower scale.
A source-selecting separation must shrink and must still dominate the
source error in the relative sense above. The existence of a growing
near-zero block does not logically exclude such a hierarchy. Neither
that hierarchy nor its required rate is in the archive.

\paragraph{Identification of the limit.}
By \eqref{eq:v150-mellin}, with
$\Xi(z)=\xi(1/2+iz)$ and the even functional equation,
\[
 \widehat\phi(z)=\frac{2}{\sqrt{6\pi}\,\|h\|_2}\,\Xi(z).
\]
In particular $\widehat\phi(0)>0$. Once either selection criterion is
proved, $\widehat v_a(0)\ne0$ eventually and
\[
 \frac{\widehat v_a(z)}{\widehat v_a(0)}
 \longrightarrow\frac{\Xi(z)}{\Xi(0)}
 \quad\text{locally uniformly on }\mathbb C.
\]
The scalar is fixed by the stated $L^2$ normalization and then by the
value at zero; phase/scale and translations must not be left free.
If the real-zero-ground-transform theorem's complete hypotheses hold
cofinally as well, this convergence rules out a nonreal zero of $\Xi$:
a small closed disk about such a zero, disjoint from the real axis and
with zero-free boundary, would by uniform convergence and the argument
principle eventually contain the same positive number of zeros of
$\widehat v_a$, a contradiction. This is a conditional implication,
not a new RH proof. Cofinal simplicity/evenness and applicability of
that theorem are themselves additional to the single $\lambda=4$ result.

\subsection*{An exact abstract nonselection model}

The next construction tests precisely the proposed \emph{abstract}
premises: complete compact-resolvent operators on increasing intervals;
simple even nonnegative ground functions; entire ground transforms with only real
zeros; monotone positive ground energies tending to zero; a fixed
candidate with vanishing full residual; dense fixed tests with vanishing
residual; ordinary strong-resolvent collapse; and even uniformly bounded
support of the ground functions. It does \emph{not} impose the
support-consistent arithmetic Weil form, the prime identities, the
CvS distribution representation, or a positivity-improving semigroup.
Its ground functions are nonnegative, not strictly positive on all of
the increasing interval. It cannot disprove an arithmetic
selection theorem using those additional structures.

\begin{proposition}[Abstract nonselection despite fast absolute residuals]
\label{prop:ns43-ccm-nonselection}
There is a family satisfying all the abstract properties just listed,
including uniformly tight ground supports and the explicit Gaussian
cutoff candidate, whose phase-fixed normalized ground transforms have
two different subsequential limits. The absolute residual rate can be
made arbitrarily fast by a positive scalar rescaling.
\end{proposition}
\begin{proof}
Let $b=1_{[-1/4,1/4]}$ and let $g_1=b*b*b*b/\|b*b*b*b\|_2$.
Then $g_1$ is nonnegative and even, lies in $H^2(\mathbb R)$, and is
supported in $[-1,1]$. Put $g_2(x)=2^{-1/2}g_1(x/2)$, supported in
$[-2,2]$. Both are unit vectors, distinct and not scalar multiples.
Their entire transforms are nonzero constants times
$\operatorname{sinc}(z/4)^4$ and $\operatorname{sinc}(z/2)^4$,
respectively, where $\operatorname{sinc}z=\sin z/z$. Every zero is real.
Their normalized transforms have distinct zero sets.

For integers $n\ge3$, let $D_n=-d^2/dx^2$ be the Dirichlet Laplacian
on $H_n=L^2(-n,n)$, with domain $H^2(-n,n)\cap H_0^1(-n,n)$.
Its unit first eigenvector is
$h_n(x)=n^{-1/2}\cos(\pi x/(2n))$.
Let $j(n)=1$ for even $n$ and $2$ for odd $n$. Define
\[
 w_n=\frac{h_n-g_{j(n)}}{\|h_n-g_{j(n)}\|},\qquad
 U_n=I-2|w_n\rangle\langle w_n|,\qquad
 A_n=n^{-1}U_n(I+D_n)U_n.
\]
The denominator is nonzero; its square is
$2-2\langle h_n,g_{j(n)}\rangle\to2$.
Thus $U_n$ is an even-parity-preserving unitary involution with
$U_nh_n=g_{j(n)}$. Because $w_n\in\mathcal D(D_n)$, it preserves
that domain. The complete $A_n$ is positive, self-adjoint with compact
resolvent, and its simple even ground is exactly $g_{j(n)}$ with
\[
 \mu_n=\frac1n\left(1+\frac{\pi^2}{4n^2}\right)\downarrow0,
 \qquad \lambda_{2,n}-\mu_n=\frac{3\pi^2}{4n^3}>0.
\]
All ground transforms have only real zeros, but the normalized ground
functions and their normalized entire transforms alternate between two
different limits. Positivity and a fixed positive phase do not repair this.

For each fixed $f\in C_c^\infty(\mathbb R)$ and all sufficiently large
$n$, $f\in\mathcal D(D_n)$. Moreover
$\|(I+D_n)w_n\|$ is uniformly bounded: the $g_j$ have fixed finite
$H^2$ norms, while $\|D_nh_n\|=\pi^2/(4n^2)$ and the denominators
are bounded away from zero. Consequently
\[
 \|A_nf\|\le\frac1n\left(\|f-f''\|
        +2\|f\|\,\|(I+D_n)w_n\|\right)\longrightarrow0.
\]
The same estimate holds for the fixed candidate $g_1$. It also holds
uniformly for the explicit Gaussian cutoff candidate $p_n$ used above:
its $H^2$ norms are uniformly bounded, its cutoff normalization stays
away from zero, and $\|p_n\|=1$. Thus $\|A_np_n\|=O(1/n)$ while
$p_n\to\phi$ in every exponentially weighted $L^1$ norm. Even the
known explicit Gaussian candidate and its transform identity do not
select the alternating grounds of this abstract operator family.
Let $\widetilde A_n=A_n\oplus0$ on $L^2(\mathbb R)$ by zero extension.
For every nonreal $z$ and every such fixed $f$,
\[
 \|(\widetilde A_n-z)^{-1}f+z^{-1}f\|
 \le\frac{\|\widetilde A_nf\|}{|z||\Im z|}\to0.
\]
Density of $C_c^\infty$ and the uniform self-adjoint resolvent bound
extend this to all of $H$. Thus the ordinary strong-resolvent limit
is zero. This model meets every listed abstract premise, including
weighted tightness of the grounds, yet supplies no unique selection.
Its ground multiplicity is one at every $n$; the failure is not caused
by a degenerate eigenspace at a fixed window.
The scalar $1/n$ can be replaced throughout by any positive sequence
$\theta_n\downarrow0$. The eigenvalues and residuals are rescaled,
the ground functions are unchanged, and the residual estimate becomes
$O(\theta_n)$. Thus even an arbitrarily prescribed rapid absolute
residual rate cannot replace the missing relative estimate. This
rescaling does not restore arithmetic support consistency.
\end{proof}


\paragraph{Ordinary norm is also the wrong transform topology by itself.}
Even if a selection estimate gives $f_a\to\phi$ in $L^2$, locally
uniform entire convergence still needs a rate or a tail condition.
For any fixed $R>0$, add to the positive cutoff source a positive even
unit bump supported near $\pm(a-1)$ with coefficient $e^{-Ra/2}$,
and renormalize. The resulting compactly supported unit functions
converge to $\phi$ in $L^2$, while their transforms at $iR$ grow like
a positive constant times $e^{Ra/2}$. This is a separate topology
counterexample; it does not claim the real-zero property or the
arithmetic ground equation.

\subsection*{A restriction on renormalized profile selection}

A possible alternative is to renormalize the shrinking spectrum and
prove a nonzero profile form with a unique selected direction.
The ordinary zero-resolvent limit cannot provide this profile.
\emph{Scope reset for this subsection only:} $q_a$ now denotes the
\emph{full both-parity} physical Weil form on $L^2(I_a)$, not merely
the even-sector form used for the preceding selection statements.
Here $J_a$ is zero extension into the full $H=L^2(\mathbb R)$;
$\alpha_a$ remains the Rayleigh value of the even cutoff $p_a$.
Thus the use of all real translates and their density in full $H$
is exactly the setting of \ref{lem:v136-total-radicals}.
Here is a precise obstruction to choosing too coarse a scale.
\begin{proposition}[A coarse rescaled liminf cannot select one radical]
\label{prop:ns43-ccm-coarse-profile}
Let $s_a>0$ and suppose
\[
 \varepsilon_a/s_a\to0,\quad |\alpha_a|/s_a\to0,
 \qquad \varepsilon_a=Ca e^{-e^a/100},
\]
where each fixed finite translate combination may have its own constant
$C$, as in v1.35--v1.36. Suppose a nonnegative lower-semicontinuous
quadratic form $\mathcal E$ on $H$ satisfies the proposed strong
liminf inequality
\[
 \mathcal E(f)\le\liminf_a
 \frac{q_a[f_a]-\alpha_a\|f_a\|^2}{s_a}
 \quad\text{whenever }J_af_a\to f.
 \tag{S5}
\]
Then $\mathcal E\equiv0$ on $H$.
\end{proposition}
\begin{proof}
For every fixed $f$ in the dense translated-radical span, insert its
smooth cutoff $f_a$. The complete residual implies the right side is
zero, so $\mathcal E(f)=0$. Lower semicontinuity and density then imply
$\mathcal E\equiv0$ on $H$. Therefore no such coarse rescaled limit
can have one-dimensional kernel $\mathbb C\phi$.
\end{proof}
This is a scoped implication about (S5), not a prohibition on finer
scales, different compactness mechanisms, or arithmetic ground selection.

Finally, in the actual support-consistent Weil family, write
$\mu_a^{\mathrm{full}}=\inf\sigma(\mathsf W_{e^a})
=\min(\mu_a^{\mathrm{even}},\mu_a^{\mathrm{odd}})$ on the full
$L^2(I_a)$. A uniform lower bound on this \emph{complete, both-parity}
ground energy is already the bounded-floor
target of v1.36, and $\mu_a^{\mathrm{full}}\to0$ is equivalent to
full both-parity positivity
via monotonicity and the known source upper bound. Those must not be
rebranded as weak new selection assumptions. Neither conclusion alone
has been shown here to identify the Gaussian ground profile. The
abstract countermodel establishes no arithmetic nonimplication.

